\documentclass[11pt,a4paper]{article}

\usepackage[T1]{fontenc}
\usepackage[utf8]{inputenc}
\usepackage{lmodern}
\usepackage[margin=2.6cm]{geometry}
\usepackage{amsmath,amssymb,bm}
\usepackage{booktabs}
\usepackage{graphicx}
\usepackage{listings}
\usepackage{xcolor}
\usepackage{microtype}
\usepackage[hidelinks]{hyperref}
\usepackage{caption}

\definecolor{codegray}{gray}{0.96}
\definecolor{keywordblue}{rgb}{0.1,0.2,0.6}
\definecolor{commentgreen}{rgb}{0.15,0.45,0.25}

\lstset{
  language=[90]Fortran,
  basicstyle=\ttfamily\footnotesize,
  keywordstyle=\color{keywordblue}\bfseries,
  commentstyle=\color{commentgreen}\itshape,
  backgroundcolor=\color{codegray},
  frame=single,
  rulecolor=\color{gray!40},
  breaklines=true,
  columns=fullflexible,
  showstringspaces=false,
  captionpos=b,
  xleftmargin=0.5em,
  framexleftmargin=0.5em,
}

\newcommand{\zj}{z_j}
\newcommand{\jsat}{j_{\mathrm{sat}}}
\newcommand{\zjsat}{\zj^{\mathrm{sat}}}
\newcommand{\zjsh}{\zj^{\mathrm{sh}}}
\newcommand{\dd}{\mathrm{d}}
\newcommand{\Gam}{\Gamma_{\mathrm{sh}}}
\newcommand{\code}[1]{\texttt{\small #1}}

\title{\bf A Charge-Conserving Sheath Boundary Condition for JOREK\\[2mm]
\large Weak Galerkin Imposition of the Langmuir Characteristic\\
on the Boundary Trace Space of \texttt{model600}}
\author{M\'at\'e Sz\H{u}cs}
\date{26 August 2026}

\begin{document}
\maketitle

\begin{abstract}
\noindent
The scrape-off layer (SOL) electric field is set by the Debye sheath at the divertor targets, and
the resulting $\bm{E}\times\bm{B}$ drift through the private flux region is a leading candidate
mechanism for the in--out divertor density asymmetry (high-field-side high-density, HFSHD). JOREK's
reduced-MHD \code{model600} previously carried a Dirichlet condition on the electrostatic potential
at the wall, which forbids exactly this physics. This paper documents a boundary condition that
instead closes charge continuity at the wall through the Langmuir current--voltage characteristic,
allowing the wall potential to be determined self-consistently by the plasma.

The central difficulty is not the physics but the discretisation. Imposing the characteristic
\emph{pointwise at boundary nodes} was found to satisfy it there to a relative accuracy of
$1.8\times10^{-4}$ while still leaving a $33$--$48\,\%$ discrepancy between the current the sheath
passes and the current the plasma delivers, because in a $C^1$ bicubic finite-element basis the
currents are integrated \emph{between} the nodes. The resolution is to impose the characteristic
where the currents are measured: as a Galerkin projection onto the boundary trace space, replacing
the boundary rows of the current-definition equation. Three ingredients proved individually
necessary --- row \emph{replacement} rather than penalisation, row normalisation by the trace mass
diagonal (whose entries span $2.6\times 10^{10}$ on this mesh), and a trust region applied to the
residual but deliberately \emph{not} to the Jacobian.

The resulting condition closes the wall current budget to four significant figures on both targets
($I_{\mathrm{sheath}}=I_{\mathrm{Amp\grave{e}re}}$ to $0.00\,\%$), reduces the weak residual to
$6.5\times10^{-4}$ of the ion saturation current, requires no tuning parameter, and has run
$\sim\!3900$ time steps to wall-clock timeout without incident on an ASDEX~Upgrade X-point geometry.
\end{abstract}

\tableofcontents

\section{Introduction and motivation}

\subsection{Why the wall potential must be free}

In a diverted tokamak the plasma terminates on material surfaces. Because electrons are far more
mobile than ions, a surface immersed in plasma charges negatively until the electron flux is
retarded enough to match the ion flux; the resulting non-neutral layer is the Debye sheath, and the
potential drop across it is of order $3\,kT_e/e$. The sheath is thinner than any mesh cell in an MHD
code by many orders of magnitude, so it cannot be resolved --- it must be represented as a boundary
condition relating the current through the surface to the potential at the surface.

That relation matters for more than the current budget. The floating potential
$\Phi_f = \Lambda kT_e/e$ is proportional to the local electron temperature, so a target with a
radially varying $T_e$ profile carries a radially varying potential, hence a radial electric field
$E_r$ at the target, hence an $\bm{E}\times\bm{B}$ drift. In a divertor this drift runs through the
private flux region (PFR) and transports particles between the inner and outer targets. It is one of
the standard explanations for the observed high-field-side high-density (HFSHD) asymmetry, and it is
the physics this work exists to enable.

A Dirichlet condition $u = 0$ on the wall states $\Phi = 0$ everywhere on it. That is not a gauge
choice once a sheath is present --- it asserts a specific, and wrong, physical state (deep electron
saturation, $X = -\Lambda$), it eliminates $E_r$ at the target by construction, and it therefore
makes the HFSHD drift mechanism structurally inaccessible. Replacing it is a prerequisite for any
study of divertor asymmetries, detachment onset, or target biasing in this model.

\subsection{Scope}

This paper documents the implementation for \code{model600} (reduced MHD with separate ion and
electron temperatures, optional neutrals and kinetic impurities) at toroidal mode number
$n_{\mathrm{tor}} = 1$, i.e.\ axisymmetric. Section~\ref{sec:physics} sets out the sheath physics and
its expression in JOREK's normalised variables. Section~\ref{sec:bvp} identifies which discrete
equation the condition belongs in --- a question with a non-obvious answer that cost several
failed attempts. Section~\ref{sec:numerics} is the core: why the natural pointwise implementation
fails despite appearing to succeed, and how a weak formulation repairs it.
Section~\ref{sec:implementation} describes the code. Section~\ref{sec:results} reports verification.

\section{Sheath physics}
\label{sec:physics}

\subsection{The current--voltage characteristic}

For a Maxwellian plasma at a wall, the ion current is saturated at the Bohm value while the electron
current is retarded by the sheath potential according to a Boltzmann factor. The net current density
along the field, referenced to the wall, is the Langmuir probe characteristic
(Stangeby~\cite{stangeby}, eq.~2.68):
\begin{equation}
  j_\parallel = \jsat\left(1 - e^{-X}\right),
  \qquad
  X \equiv \frac{e\Phi}{kT_e} - \Lambda,
  \qquad
  \Phi = V_{\mathrm{sheath}} - V_{\mathrm{wall}} .
  \label{eq:char}
\end{equation}
Here $\Lambda$ is the floating-potential coefficient. Its Maxwellian value is
\begin{equation}
  \Lambda_0 = \ln\sqrt{\frac{m_i}{2\pi m_e}} \;\simeq\; 3.19 \quad\text{(deuterium)} ,
  \label{eq:lambda0}
\end{equation}
and an optional local correction accounting for the ion contribution to the sound speed is available:
\begin{equation}
  \Lambda = \Lambda_0 - \tfrac12\ln\!\left(\frac{\gamma\,(T_i+T_e)}{T_e}\right),
  \qquad
  \frac{\partial\Lambda}{\partial T_i} = -\frac{1}{2(T_i+T_e)},
  \qquad
  \frac{\partial\Lambda}{\partial T_e} = \frac{T_i}{2T_e(T_i+T_e)} .
\end{equation}
Equation~\eqref{eq:char} has $j_\parallel = 0$ at $X=0$, i.e.\ $\Phi = \Lambda kT_e/e$: the floating
potential. At $\Phi = 0$ (a grounded wall, $X = -\Lambda$) the sheath is in electron saturation and
passes $j_\parallel = \jsat(1-e^{\Lambda}) \approx -19\,\jsat$.

The ion saturation current density follows the Bohm condition,
\begin{equation}
  \jsat = e\, n\, c_s\, g(b_n), \qquad c_s = \sqrt{\gamma\,(T_i+T_e)/m_i},
\end{equation}
where $g(b_n)$ is the Chodura--Riemann obliqueness function --- the same factor the Mach-1 boundary
condition uses for the parallel outflow velocity --- carrying the sign of $\bm{B}\!\cdot\!\bm{n}$.
Using the same $g$ in both places is what makes the current and the particle flux mutually
consistent at the same surface, and it is why \code{bcs\%mach1} must be enabled on any boundary type
carrying this condition.

\subsection{Forward versus inverse form: a decisive choice}
\label{sec:forward}

Equation~\eqref{eq:char} may be used in either direction. The \emph{inverse} form, solving for the
potential given the current,
\begin{equation}
  \frac{e\Phi}{kT_e} = \Lambda - \ln\!\left(1 - \frac{j_\parallel}{\jsat}\right),
\end{equation}
is the more intuitive statement --- ``the wall floats to whatever potential passes the delivered
current'' --- and it was how the first implementation was written. It is also singular precisely
where divertor plasmas live: as $j_\parallel \to \jsat$ the logarithm diverges, and for
$j_\parallel > \jsat$ there is \emph{no solution at all}. A restart equilibrium constructed without
this boundary condition has no reason to respect $|j_\parallel| \le \jsat$, and in practice does not.

The forward form used here, $j_\parallel = j_\parallel(\Phi)$, is single-valued for all $\Phi$, and
its derivative
\begin{equation}
  \frac{\partial j_\parallel}{\partial \Phi} \propto e^{-X}
\end{equation}
tends smoothly to zero on the ion side. No clipping is needed there, and the linearisation never
becomes singular. This choice is the reason the boundary condition is robust from an arbitrary
restart state, and it is worth stating explicitly because the two forms are mathematically
equivalent and numerically utterly different.

\subsection{Finite conductance at ion saturation}

The forward characteristic saturates \emph{exactly} at $\jsat$: no potential, however large, passes
more ion current. A plasma delivering even a few percent more than $\jsat$ therefore has no
consistent boundary state, and the residual settles at a small constant that drives $u$ linearly
forever. This is not hypothetical --- it is the observed failure mode of the first working
configuration.

The remedy is to give the ion branch a finite differential conductance by adding a softplus term,
\begin{equation}
  f(X) = \left(1 - e^{-X}\right) + s\,\ln\!\left(1 + e^{X}\right),
  \qquad
  \zjsh = \zjsat\, f(X),
  \label{eq:fchar}
\end{equation}
with $s$ = \code{sheath\_sat\_slope}. Because $\ln(1+e^X)$ is exponentially small for $X \ll 0$, the
electron branch and the floating potential are essentially untouched, while $f$ becomes unbounded
above so that every demanded current is reachable at a finite $X$. The function is $C^\infty$ and is
applied to $f$ and $\dd f/\dd X$ together.

On the electron side the exponential is instead bounded by a smooth limiter,
\begin{equation}
  X_{\mathrm{lim}} = X_{\min} + \delta X \ln\!\left(1 + e^{(X - X_{\min})/\delta X}\right),
\end{equation}
which is $C^1$, monotone, has $\dd X_{\mathrm{lim}}/\dd X \in (0,1)$, and tends to the identity for
$X \gg X_{\min}$. The ion side needs no such limiter, which is the whole point of
Section~\ref{sec:forward}.

Two numerical details are worth recording. First, $1-e^{-X}$ cancels catastrophically near $X=0$ ---
which is the floating condition, i.e.\ the regime of interest rather than a corner case. Fortran has
no \code{expm1} intrinsic, so a four-term Taylor series is used below $|X| < 10^{-4}$, giving a
relative error of $\sim\!10^{-18}$. Second, all exponentials are evaluated only within
$|X| \le 30$; outside, the asymptotic branch is taken, which keeps a debug build with
\code{-ffpe-trap=overflow,denormal} alive for any state the nonlinear solver passes through.

\subsection{JOREK normalisation and the sign convention}

JOREK's reduced-MHD variables are normalised; the sheath must be expressed in them. The
electrostatic potential enters through the stream function $u$. The JOREK reference paper
(H\"olzl \emph{et al.}~\cite{hoelzl2021}, eq.~26) defines $u = \Phi/F_0$ with a velocity ansatz
$\bm{v}_{\mathrm{pol}} = -R\,\nabla u \times \bm{e}_\phi$, whereas the \code{model600} element matrix
implements $\bm{v}_{\mathrm{pol}} = +R\,\nabla u \times \bm{e}_\phi$. The code's $u$ is therefore
\emph{minus} the paper's, and
\begin{equation}
  \boxed{\;\Phi = -F_0\,u\;}
\end{equation}
The same sign inversion applies to $\psi$ and to the current variable: the code's
$\zj = \Delta^{*}\psi$ is minus the paper's $j$. Getting this wrong inverts the entire boundary
condition, so it is stated here explicitly and derived once, in one module, used by all three
consumers (assembly, diagnostic, and the legacy nodal route) so that they cannot drift apart.

With $n_0$ the central density, $\rho_0 = n_0 m_i$, and $m_i = A\,u_{\mathrm{amu}}$:
\begin{align}
  \frac{e\Phi}{kT_e} &= \frac{\tfrac12 a_n u - v_w}{T_e},
  &
  a_n &= -\,\frac{2 e F_0 \sqrt{\mu_0 \rho_0}}{m_i},
  \label{eq:an}\\[1mm]
  v_w &= e\,V_{\mathrm{wall}}\,\mu_0\,n_0,
  &
  c_{\mathrm{sat}} &= -\tfrac12 a_n,
\end{align}
and the saturation current in code variables (Artola \emph{et al.}~\cite{artola}, eqs.~5 and 16,
in the code's sign) is
\begin{equation}
  \zjsat = c_{\mathrm{sat}}\,\rho\,g(b_n)\,\frac{c_s}{|\bm{B}|},
  \qquad c_s = \sqrt{\gamma\,(T_i + T_e)} .
  \label{eq:zjsat}
\end{equation}
This corresponds to a force-free, field-aligned SOL current, $J_\parallel = -j\bm{B}/F_0$ in the
paper's variables, which is the physically appropriate interpretation in the scrape-off layer.

\subsection{Temperature floors}

The characteristic divides by $T_e$ and is therefore sensitive to the positivity floor applied to
the temperature. JOREK's global floor \code{T\_min\_neg} exists to keep every other term
well-behaved and is normally set well \emph{above} divertor temperatures: with
\code{T\_min\_neg}~$=3\times10^{-5}$ the knee sits at $1.93$~eV, so a real target profile running
$0.3 \to 1.5$~eV is delivered to the sheath as $1.14 \to 1.58$~eV. Since $\Phi_f = \Lambda kT_e/e$
inherits that compression, the target's \emph{radial potential gradient} --- precisely the $E_r$ that
drives the PFR drift --- is flattened by a factor $\sim\!3.5$ before the boundary condition ever sees
it. A separate floor \code{T\_min\_sheath}, with the same functional form, is therefore provided so
that the sheath can see a colder plasma than the rest of the model. This is a physics-relevant knob,
not a numerical convenience.

\section{Where the condition belongs in the discrete system}
\label{sec:bvp}

\subsection{Two candidate equations}

The sheath supplies one scalar relation per boundary point. It can be attached either to the
vorticity (charge continuity) equation, as a surface flux, or to the current-definition equation
$\zj = \Delta^{*}\psi$, as an algebraic constraint. Both were implemented; the distinction is not
cosmetic.

In \code{model600} the vorticity equation is assembled in \emph{strong} form. It therefore has no
boundary flux for a surface term to replace, and adding one injects a spurious source at the wall
rather than supplying a missing one. This route (\code{natural\%u}) works, but it is a Robin-type
penalty whose strength must be capped against the row's own polarisation diagonal to prevent $u$
from being slaved pointwise.

The route pursued here attaches the characteristic to the current definition instead: the boundary
rows of the $\zj$ equation are replaced by the statement $\zj = \zjsh$. This is an algebraic
constraint --- no $\theta$-scheme, no time step --- and it puts the condition on the variable the
characteristic actually predicts.

\subsection{Why \texttt{natural\%zj} is required: the open loop}

A subtlety here falsified an earlier assumption and is worth recording. It was originally believed
that \code{dirichlet\%zj} could remain enabled, on the grounds that the frozen $\zj$ trace cancels
exactly between the strong-form volume term and the added surface flux. This is \emph{false}, and
the diagnostic falsified it directly: with \code{dirichlet\%zj = .true.}, the Amp\`ere current
$I_{\mathrm{Amp}} = \oint \zj\,(\bm{B}\!\cdot\!\bm{n})\,\dd S$ stayed constant to four digits for an
entire run while the sheath current ran from $-15540$~A to $-23$~A. The sheath adapted all the way to
floating, the plasma current could not move at all, and $u$ diverged trying to reconcile them.

The $j$--$V$ loop is \emph{open} whenever $\zj$ at the wall is pinned: the boundary condition can
then only move the potential, and if the delivered current exceeds $\jsat$ no potential satisfies the
characteristic. The current-definition surface term must therefore be restored
(\code{natural\%zj = .true.}, \code{dirichlet\%zj = .false.}), coupling $\zj$ at the boundary to
$\nabla\psi\!\cdot\!\bm{n}$ --- the \emph{normal} derivative of $\psi$, which \code{dirichlet\%psi}
does not pin, since that fixes the value and the tangential derivative only. That normal derivative
is the degree of freedom through which the wall current responds to the sheath.

\subsection{Required namelist configuration}

\begin{lstlisting}[caption={Enabling the weak sheath condition on the boundary types carrying the strike points.}]
  bcs(1)%sheath_zj_weak = .true.    ! the weak Galerkin route (this work)
  bcs(1)%natural%zj     = .true.    ! REQUIRED - without it the j-V loop is open
  bcs(1)%dirichlet%zj   = .false.   ! ... and its Dirichlet has to come off
  bcs(1)%dirichlet%u    = .false.   ! u is free, set by charge continuity at the wall
  bcs(1)%dirichlet%w    = .true.    ! KEEP; leave natural%w = .false.
  bcs(1)%mach1          = .true.    ! j_sat assumes Mach 1 at the same wall (default for type 1)
  bc_natural_open       = .true.    ! the boundary integrals live in that branch
\end{lstlisting}

At least one boundary type must retain \code{dirichlet\%u = .true.} --- typically the main chamber
wall. The reason is that $u$ enters the vorticity equation only through its gradient, and the sheath
term loses its grip on $u$ in ion saturation; without a gauge somewhere, the constant mode of $u$ is
undetermined.

\section{Numerics: from pointwise to weak}
\label{sec:numerics}

\subsection{The discretisation and its trace space}

JOREK uses $C^1$-continuous bicubic B\'ezier finite elements. Each node carries four degrees of
freedom per variable --- value, $\partial_s$, $\partial_t$, and the mixed $\partial_s\partial_t$ ---
so a boundary edge of an element sees, in principle, all of them.

The key structural fact is that on a boundary edge the \emph{trace} of a bicubic is determined by
the value and \emph{tangential}-derivative degrees of freedom at the two endpoints only. The
normal-derivative and mixed degrees of freedom have vanishing basis functions on the edge. The trace
space of a boundary segment is therefore spanned by four functions per element edge: two nodes
$\times$ two trace DOFs. This is what makes a boundary projection tractable, and it defines the
column set of everything that follows.

\subsection{The pointwise route, and why it deceives}
\label{sec:nodalfail}

The natural implementation imposes $\zj = \zjsh$ at each boundary \emph{node}, as a linearised
constraint row
\begin{equation}
  \dd \zj - \sigma \sum_k \frac{\partial \zjsh}{\partial x_k}\,\dd x_k = \sigma\,\left(\zjsh - \zj\right),
\end{equation}
with the diagonal left unscaled so that $\sigma$ is the per-step closure fraction.

Measured, this works extremely well \emph{on its own terms}. The pointwise residual
\begin{equation}
  \varepsilon_{\mathrm{node}}
  = \sqrt{\frac{\sum_{\mathrm{nodes}} (\zj - \zjsh)^2}{\sum_{\mathrm{nodes}} (\zjsat)^2}}
\end{equation}
falls to $1.8\times10^{-4}$. And yet the physical acceptance test --- that the current the sheath
passes equals the current the plasma delivers,
\begin{equation}
  I_{\mathrm{sheath}} = \oint_{\Gam} \frac{\zjsh\,(\bm{B}\!\cdot\!\bm{n})}{F_0\mu_0}\,\dd S
  \;\stackrel{!}{=}\;
  I_{\mathrm{Amp}} = \oint_{\Gam} \frac{\zj\,(\bm{B}\!\cdot\!\bm{n})}{F_0\mu_0}\,\dd S
  \label{eq:closure}
\end{equation}
--- failed by $33$--$48\,\%$ on the inner target, run after run, under every variation of the
characteristic, the gates and the relaxation.

The explanation is that the two statements are simply different. The constraint is imposed at nodes;
the currents in~\eqref{eq:closure} are integrated at Gauss points \emph{between} the nodes, where the
cubic trace is unconstrained. A cubic passing exactly through the right values at its endpoints can
still integrate to something quite different. The corresponding Gauss-point residual
$\varepsilon_{\mathrm{gauss}}$ (same definition, weighted by $\dd S$) was measured at $\approx 2$ ---
four orders of magnitude larger than the nodal one.

This was confirmed decisively by instrumenting the \emph{weak} residual as a pure diagnostic before
changing anything: as $\varepsilon_{\mathrm{node}}$ fell $1.11 \to 0.08 \to 1.8\times10^{-4}$, the
Galerkin residual stayed at $\mathcal{O}(1)$. The constraint was being satisfied where it was
imposed and nowhere else.

\subsection{The weak formulation}

Let $\{N_a\}$ span the boundary trace space and $\Gam$ be the sheath boundary. Define the Galerkin
residual and the trace mass diagonal
\begin{equation}
  F_a = \int_{\Gam} N_a \left(\zjsh - \zj\right) \dd S,
  \qquad
  D_a = \int_{\Gam} N_a N_a \,\dd S .
  \label{eq:FaDa}
\end{equation}
The condition imposed is $F_a = 0$ for every trace DOF $a$ --- the orthogonal projection of the
sheath mismatch onto the trace space. Because this is enforced exactly where the currents are
integrated, the closure~\eqref{eq:closure} follows by construction rather than by convergence.

Linearising about the current state, the replacement row for trace DOF $a$ is
\begin{equation}
  \underbrace{\int_{\Gam} N_a N_b \,\dd S}_{\text{$\zj$ column}} \dd z_{j,b}
  \;-\;
  \sum_{x \in \{u,\rho,T_i,T_e\}} \underbrace{\int_{\Gam} N_a \frac{\partial \zjsh}{\partial x} N_b \,\dd S}_{\text{state columns}} \dd x_b
  \;=\; F_a .
  \label{eq:row}
\end{equation}
Note that no $\theta$ and no time step appear: this is an algebraic constraint on the state, not a
discretised evolution equation.

\subsection{Ingredient 1: replacement, not penalisation}

The obvious way to impose~\eqref{eq:row} is as a penalty, adding $\beta F_a$ to the existing $\zj$
equation. This cannot work here, and understanding why is the crux.

The volume equation $\zj = \Delta^{*}\psi$ is integrated by parts, and its boundary surface term is
\emph{refused} by the assembly: the term's Jacobian requires columns at the normal-derivative DOFs,
which the trial-function loop cannot produce (it emits columns only at the tangential DOFs). This is
harmless while \code{dirichlet\%zj} freezes the trace, because those rows are overwritten anyway ---
and the code says exactly this in its comments. But the weak route \emph{must} release that
Dirichlet (Section~\ref{sec:bvp}), and the boundary equation is then \emph{incomplete}. Adding a
penalty to an incomplete equation cannot repair it. Measured: $\beta=10^{-2}$ left the wrong
equation dominant; $\beta=1$ produced a period-2 divergence.

Writing the row with the code's \code{zbig}~$=10^{12}$ multiplier annihilates the incomplete
equation instead of adding to it --- exactly as every Dirichlet condition in JOREK already does.

It is worth noting explicitly that JOREK's ``Dirichlet'' is itself a penalty: \code{mod\_assembly}
sets a single matrix entry to \code{zbig} rather than clearing the row. The remaining entries survive
but are $10^{18}$ times smaller. The same mechanism is reused here, applied to a full row rather than
one entry.

\subsection{Ingredient 2: row normalisation}
\label{sec:normalisation}

A Galerkin trace block has internal scale structure that a pointwise Dirichlet does not. Its mass
matrix entries scale as
\begin{equation}
  \int N_a N_b \,\dd S \sim
  \begin{cases}
    h & \text{value} \times \text{value},\\
    h^2 & \text{value} \times \text{derivative},\\
    h^3 & \text{derivative} \times \text{derivative},
  \end{cases}
\end{equation}
so at $h \sim 1$~mm the block spans several orders of magnitude \emph{internally}. The measured
spread on this mesh is
\begin{equation}
  \frac{D_{\max}}{D_{\min}} = \frac{1.62\times10^{-3}}{6.34\times10^{-14}} = 2.6\times10^{10},
\end{equation}
four orders beyond the $10^6$ estimated from the $h$-scaling alone (the true scaling involves
\code{element\_size}$^2$, not $h^2$). No uniform coefficient can work across such a range:
$\beta = 10^{9}$ died in a single step.

Dividing each row by its own diagonal,
\begin{equation}
  \frac{1}{D_a}\left[\int N_a N_b \dd S\,\dd z_{j,b} - \sum_x \int N_a \frac{\partial \zjsh}{\partial x} N_b \dd S\,\dd x_b\right]
  = \frac{F_a}{D_a},
\end{equation}
makes every row $\mathcal{O}(1)$, after which a single \code{zbig} makes them all uniformly dominant
and \emph{no penalty parameter is required}. This is the reason the final scheme is tuning-free.

A degeneracy floor at $10^{-14}D_{\max}$ guards against a row that received essentially no boundary
support, where $1/D_a$ would manufacture an enormous row out of numerical noise. In production runs
this floor is never triggered: the diagnostic reports \code{204 accumulated, 204 replaced, 0 below
the floor} at every step.

\subsection{Ingredient 3: a trust region on the residual only}

The residual $\zjsh - \zj$ is unbounded on the electron branch. At $u = 0$ the characteristic sits at
$X = -\Lambda$ with $f = 1 - e^{\Lambda} \approx -19$, so a restart whose wall is not already near
the floating potential asks the constraint for $\sim\!20\,\jsat$ of driving force on the very first
step. Measured doing exactly that: $e\Phi/kT_e$ collapsed to $0.00/0.01/0.03$,
$100\,\%$ of the area electron-limited, $I_{\mathrm{sheath}} = -12$~kA against
$I_{\mathrm{Amp}} = +1$~kA, and blow-up in four steps.

The residual is therefore passed through
\begin{equation}
  r \;\longmapsto\; \frac{r}{1 + |r| / (r_{\max}\,\zjsat)},
  \label{eq:trust}
\end{equation}
with $r_{\max}$ = \code{sheath\_weak\_rmax} (default 2). This map is the \emph{identity} for
$|r| \ll r_{\max}\zjsat$, so the fixed point and the local convergence rate are untouched, and it
saturates at $r_{\max}\zjsat$ however far off the characteristic the state starts. An algebraic
rather than a $\tanh$ clip is used deliberately: the derivative of~\eqref{eq:trust} decays only
algebraically, whereas a $\tanh$'s exponential decay would leave the row effectively explicit during
the approach.

\paragraph{The Jacobian is deliberately \emph{not} bounded.} This is a trap worth flagging. The
Newton step solves $J\,\dd x = F$; damping $J$ by a factor $\alpha$ \emph{magnifies} $\dd x$ by
$1/\alpha$, which is the opposite of what a trust region is for. Bounding $F$ alone caps the step at
$r_{\max}\jsat$ and is exact Newton near the solution.

\subsection{Ingredient 4: the validity weight}

Where the plasma demands $|\zj/\zjsat|$ far outside what the characteristic can deliver, the relation
carries no information, and driving $\zj$ toward it is actively destructive. Measured: a local
$\jsat$ collapse took the inner target from $\max|j/\jsat| = 13 \to 18.7 \to 354$ in three steps and
ended an otherwise healthy 550-step run.

The trust region cannot cover this case, because it bounds the residual by $r_{\max}z^{\mathrm{sat}}_{j,\mathrm{loc}}$,
which vanishes together with $\zjsat$ itself. A separate smooth validity weight is therefore applied,
\begin{equation}
  w = \left[1 + \left(\frac{|\zj/\zjsat|}{\mathcal{R}_{\max}}\right)^{4}\right]^{-1},
  \label{eq:weight}
\end{equation}
with $\mathcal{R}_{\max}$ = \code{sheath\_zj\_ratio\_max}. It is monotone, equals 1 to within
$6\,\%$ for ratios below $\tfrac12\mathcal{R}_{\max}$, and --- critically --- is applied to $F_a$,
$D_a$ \emph{and} the Jacobian columns alike, so the normalised row remains a consistent weighted
residual and simply fades out at pathological points, leaving those Gauss points to the assembled
equation.

\subsection{The column set: a deliberate quasi-Newton}
\label{sec:columns}

The sheath current is a function
$\zjsh(u, \rho, T_i, T_e; g_{b_n}, |\bm{B}|, R)$. The Jacobian columns emitted are therefore $\zj$
(the unit diagonal), $u$, $\rho$ and $T_i, T_e$ --- or $T$ without \code{WITH\_TiTe}. Two omissions
are worth stating:

\begin{itemize}
  \item $v_\parallel$ and $w$ appear \emph{not at all}, because \code{sheath\_current} takes neither
        as an argument; their derivatives are identically zero. This is exact, not an approximation.
  \item $\psi$ is omitted \emph{by choice}. Both $\zjsat \propto 1/|\bm{B}|$ and $g(b_n)$ do depend on
        the magnetic field, but \code{sheath\_current} returns no $\partial\zjsh/\partial\psi$, so the
        boundary geometry is frozen within a Newton iteration. This makes the scheme quasi-Newton
        rather than full Newton. It converges perfectly well in practice; it is recorded here because
        it is the most plausible source of any residual stiffness.
\end{itemize}

For a shared trace DOF the accumulator collects columns from \emph{both} adjacent edges:
3 distinct nodes $\times$ 2 trace DOFs $\times$ 5 variables = 30 columns, rising to 40 at a junction
where three boundary segments meet.

\subsection{On under-relaxation}
\label{sec:relax}

An under-relaxation parameter \code{sheath\_weak\_relax}~$=\theta$ is provided, which scales the
state columns and the residual in~\eqref{eq:row} while leaving the $\zj$ column (and hence the $D_a$
normalisation) untouched, so that the row reads
\begin{equation}
  \zj^{\,\mathrm{new}} = \zj^{\,\mathrm{old}} + \theta\left(z_{j,\mathrm{lin}}^{\mathrm{sh}} - \zj^{\,\mathrm{old}}\right).
\end{equation}
It was added to damp a period-2 step-to-step alternation observed around step 568.

\textbf{It should be left at $\theta = 1$.} The alternation proved to be a transient of the
relaxation phase, not an instability: it disappeared on its own at $\theta=1$ by step 3900. Measured
head-to-head at timeout, $\theta = 0.5$ was worse on every metric (Table~\ref{tab:relax}).
Under-relaxation makes the constraint row \emph{lag} a state that is still slowly evolving, so it
never catches up. The parameter is retained only as insurance against a genuine over-correction, and
its documentation in \code{phys\_module} records this measurement so the experiment is not repeated.

\section{Implementation}
\label{sec:implementation}

\subsection{Module map and call ordering}

\begin{table}[h]
\centering
\small
\begin{tabular}{@{}lp{9.6cm}@{}}
\toprule
\textbf{File} & \textbf{Role} \\
\midrule
\code{model600/mod\_sheath\_bc.f90}          & All sheath physics: normalisation constants, $\Lambda$, the $X$-limiter, and \code{sheath\_current}, which returns $\zjsh$, $\zjsat$, $X$ and the four derivatives. Stateless by design. \\
\code{model600/mod\_boundary\_matrix\_open.f90} & Boundary element assembly. Evaluates the characteristic at each boundary Gauss point, forms $D_a$, $F_a$ and the Jacobian blocks, and hands them to the accumulator. \\
\code{model600/mod\_sheath\_trace.f90}       & The accumulator: sums contributions across adjacent edges, normalises by $D_a$, writes the rows with \code{zbig}. \\
\code{model600/mod\_boundary\_conditions.f90}& Owns \code{a\_mat} and \code{RHS\_loc}; calls \code{sheath\_trace\_apply}. Also hosts the legacy nodal route. \\
\code{model600/mod\_sheath\_diag.f90}        & Diagnostics: $I_{\mathrm{sheath}}$, $I_{\mathrm{Amp}}$, $e\Phi/kT_e$ statistics, inner/outer split, nodal and weak residuals. \\
\code{matrix/construct\_matrix\_mod.f90}     & Resets the accumulators before the element loop; reports afterwards. \\
\bottomrule
\end{tabular}
\caption{Files involved in the weak sheath boundary condition.}
\end{table}

The ordering within one matrix construction is essential and is worth stating, because two of the
accumulators must be zeroed at different points for this reason:
\begin{enumerate}\itemsep2pt
  \item \code{sheath\_diag\_reset}, \code{sheath\_trace\_reset};
  \item the (OpenMP-threaded) element loop, which calls \code{boundary\_matrix\_open} $\rightarrow$
        \code{sheath\_trace\_add};
  \item \code{sheath\_diag\_report};
  \item \code{boundary\_conditions} $\rightarrow$ \code{sheath\_trace\_apply} $\rightarrow$
        \code{sheath\_trace\_report}.
\end{enumerate}

\subsection{The accumulator}

Rows cannot be written during the element loop, because a trace DOF receives contributions from two
adjacent boundary edges belonging to different elements, and normalisation by $D_a$ requires the
\emph{completed} sum. The accumulator is therefore keyed by global DOF index
(\code{nodes(i)\%index(direction(j))} --- the same map the nodal route uses):

\begin{lstlisting}[caption={The accumulator state (\texttt{mod\_sheath\_trace.f90}).}]
  integer, parameter :: st_max_row = 4000  ! trace DOFs per rank; 204 observed
  integer, parameter :: st_max_col = 64    ! 3 nodes x 2 trace DOFs x 5 vars = 30; junction 40

  integer, save :: st_n, st_row(st_max_row), st_nc(st_max_row)
  real*8,  save :: st_D(st_max_row), st_F(st_max_row)
  integer, save :: st_col(st_max_col, st_max_row), st_var(st_max_col, st_max_row)
  real*8,  save :: st_val(st_max_col, st_max_row)
\end{lstlisting}

Slot lookup is a backwards linear search, which is effectively $\mathcal{O}(1)$ because adjacent
boundary edges share DOFs and are visited consecutively. Overflow of either dimension is a
\emph{fatal} error rather than a silent truncation, on the grounds that a truncated row is a wrong
equation rather than a missing one.

The application step is then simply
\begin{lstlisting}[caption={Normalisation and row writing (\texttt{sheath\_trace\_apply}).}]
    sc = 1.d0 / st_D(is)

    do ic = 1, st_nc(is)
      call boundary_conditions_add_one_entry(                          &
             st_row(is), var_zj, in, st_col(ic,is), st_var(ic,is), in, &
             zbig * st_val(ic,is) * sc, index_min, index_max, a_mat)
    enddo

    call boundary_conditions_add_RHS(                                  &
             st_row(is), var_zj, in, index_min, index_max, RHS_loc,    &
             zbig * st_F(is) * sc, a_mat%i_tor_min, a_mat%i_tor_max)
\end{lstlisting}

\subsection{Thread safety}

The element loop in \code{construct\_matrix\_mod} is OpenMP-threaded. Every \code{st\_*} array is
module-level shared state that \code{sheath\_trace\_add} both searches and extends, so without
protection two threads race on \code{st\_n} and write past each other --- producing corrupted row
indices, which is a crash rather than a wrong answer. The routine body is wrapped in
\code{!\$omp critical (sheath\_trace\_accumulate)}. Boundary elements are a small fraction of the
mesh, so the serialisation cost is negligible. (\code{sheath\_diag\_add} has always been protected
the same way; the omission here cost one queue run.)

\subsection{Diagnostics}

Because the failure mode of Section~\ref{sec:nodalfail} was invisible in every quantity being
measured at the time, the diagnostic set was deliberately extended before the implementation was
changed. Per time step the code reports:

\begin{itemize}\itemsep1pt
  \item $I_{\mathrm{sheath}}$ and $I_{\mathrm{Amp}}$ per boundary type, and split by inner/outer
        target on a major-radius threshold \code{sheath\_diag\_R\_split}. The split matters because
        boundary type 1 carries \emph{both} targets, so any total averages them and opposite errors
        cancel --- useless for HFSHD, which \emph{is} an in--out difference.
  \item $e\Phi/kT_e$ min / mean / max, and separately the max \emph{where the sheath term is active},
        which distinguishes ``the plasma is overdriving the sheath'' from ``$u$ has no boundary
        condition at all here and the null space is running away''.
  \item $\max|j/\jsat|$, globally and per target: the solvability ratio. Values above 1 mean no
        potential satisfies the characteristic at that point.
  \item The fraction of wetted area sitting on the electron-side limiter.
  \item Both residuals, $\varepsilon_{\mathrm{node}}$ and $\varepsilon_{\mathrm{gauss}}$, and the weak
        residual $|F_a/D_a| / |S_a/D_a|$ with $S_a = \int N_a \zjsat \dd S$ as the reference scale.
  \item Row bookkeeping: rows accumulated, rows replaced, rows below the degeneracy floor,
        $D_{\min}$ and $D_{\max}$.
\end{itemize}

Reporting $\varepsilon_{\mathrm{node}}$ and $\varepsilon_{\mathrm{gauss}}$ \emph{separately} is what
made the diagnosis possible: node-small with Gauss-large is an interpolation problem, whereas both
large means the row is not being satisfied at all. No single pointwise number can distinguish them.

\subsection{Known limitation: MPI}
\label{sec:mpi}

Rows are written only by the rank that owns them, so a trace DOF whose two adjacent boundary edges
live on different ranks is assembled from the local edge only. Both $F_a$ and $D_a$ lose the same
contribution, so the \emph{normalised} row remains a valid weighted-residual statement --- simply
tested against a function supported on one edge rather than two. A full fix requires a halo exchange
of the accumulator.

Empirically the effect is not measurable at production sizes: runs on 16 and 32 ranks gave identical
output to four significant figures with the same 204 trace rows. This is a useful, if unplanned,
decomposition-independence check, but it is not a proof, and the limitation should be removed before
using strongly refined boundary meshes or much larger rank counts.

\section{Verification and results}
\label{sec:results}

All results are for an ASDEX Upgrade X-point geometry (shot 38773), \code{model600},
$n_{\mathrm{tor}} = 1$, boundary type 1, with kinetic impurities disabled.

\subsection{Convergence of the constraint}

\begin{table}[h]
\centering
\small
\begin{tabular}{@{}lrrrrr@{}}
\toprule
Step & $\varepsilon_{\mathrm{weak}}$ & $\varepsilon_{\mathrm{gauss}}$ & INNER closure & OUTER closure & $I_{\mathrm{wall}}$ \\
\midrule
0      & 1.759                 & 1.990                 & ---       & ---       & --- \\
9      & 0.049                 & 0.071                 & 4.1\,\%   & 0.05\,\%  & --- \\
30     & $9.9\times10^{-4}$    & $9.7\times10^{-3}$    & 0.014\,\% & 0.012\,\% & $-136$~A \\
568    & $2.0\times10^{-2}$    & $3.9\times10^{-2}$    & 0.6--1.6\,\% & 0.2--0.3\,\% & $-2950$~A \\
$\sim$3900 & $\bm{6.5\times10^{-4}}$ & $\bm{9.0\times10^{-3}}$ & $\bm{0.00\,\%}$ & $\bm{0.00\,\%}$ & $-684$~A \\
\bottomrule
\end{tabular}
\caption{Convergence of the weak sheath condition. ``Closure'' is
$|I_{\mathrm{sheath}} - I_{\mathrm{Amp}}| / |I_{\mathrm{sheath}}|$ on each target. The excursion at
step 568 is the plasma relaxing on a transport timescale, not a degradation of the constraint; it
recovers without intervention. The run terminated on wall-clock timeout, not instability.}
\label{tab:conv}
\end{table}

At the final state $I_{\mathrm{sheath}}$ and $I_{\mathrm{Amp}}$ agree to four significant figures on
both targets (inner $-281.7 / -281.7$~A, outer $-402.2 / -402.2$~A) and on the wall total
($-683.9 / -683.9$~A). For comparison, the converged \emph{nodal} route left $48\,\%$ on the inner
target and $4.1\,\%$ on the outer.

That $\varepsilon_{\mathrm{gauss}}$ collapses along with the weak residual (from 1.99 to
$9.0\times10^{-3}$) was better than expected: a Galerkin projection only has to remove the component
of the mismatch \emph{lying in} the trace space, and is under no obligation to reduce the pointwise
error. That it does so anyway indicates the trace space represents $\zjsh$ well on this mesh.

\subsection{Under-relaxation comparison}

\begin{table}[h]
\centering
\small
\begin{tabular}{@{}lrr@{}}
\toprule
Metric at timeout & $\theta = 1.0$ & $\theta = 0.5$ \\
\midrule
$\varepsilon_{\mathrm{weak}}$        & $\bm{6.5\times10^{-4}}$ & $6.3\times10^{-3}$ \\
$\varepsilon_{\mathrm{gauss}}$       & $\bm{9.0\times10^{-3}}$ & $1.5\times10^{-2}$ \\
Target closure                       & $\bm{0.00\,\%}$         & 0.05--0.12\,\% \\
Boundary type 3 closure              & $\bm{0.14\,\%}$         & 2.9\,\% \\
$e\Phi/kT_e$ max                     & $\bm{6.56}$ (falling)   & 7.32 (rising) \\
\code{construct\_matrix} time        & $\bm{0.276}$~s          & 0.334~s \\
\bottomrule
\end{tabular}
\caption{Under-relaxation of the trace row is counterproductive; see Section~\ref{sec:relax}.
Both configurations ran to timeout without crashing.}
\label{tab:relax}
\end{table}

\subsection{Physical state}

At the converged state the wall carries a net current of $-684$~A, with the inner target drawing
$-281.7$~A and the outer $-402.2$~A. The mean normalised sheath drop is $e\Phi/kT_e = 2.46$ against
$\Lambda = 3.19$, with an inner/outer contrast of $2.52$ versus $2.38$. Some $8.3\,\%$ of the wetted
area sits on the electron-side limiter, with $\max|j/\jsat| = 13.1$, steady.

Two cautions attach to these numbers. First, the large in--out asymmetries seen during the transient
(a factor 3.5 at step 30, and a reversed contrast at step 568) are \emph{not} representative of the
converged state and should not be quoted. Second, a small contrast in $e\Phi/kT_e$ does not imply a
small potential difference: $\Phi = (e\Phi/kT_e)\cdot kT_e/e$, so at steady state the asymmetry
resides in the target $T_e$ contrast rather than in the normalised drop. Assessing the HFSHD
mechanism requires the radial profile of $\Phi$ along each target --- and hence $E_r$ --- which is
not visible in any scalar reported here.

\section{Limitations and outlook}

\paragraph{Axisymmetric only.} The accumulator handles a single toroidal index, checked at setup.
Extending to $n_{\mathrm{tor}} > 1$ requires a toroidal index dimension on the accumulator and, more
substantially, thought about whether the trace projection should couple harmonics.

\paragraph{MPI halo.} See Section~\ref{sec:mpi}. Measurably harmless at 16--32 ranks; should be
fixed before boundary-refined meshes.

\paragraph{Quasi-Newton.} The omission of the $\psi$ column (Section~\ref{sec:columns}) freezes the
boundary geometry within an iteration. This is the most plausible remaining source of stiffness and
would be the first thing to add if convergence ever degrades.

\paragraph{Missing physics elsewhere in the model.} Two gaps identified in the course of this work
are independent of the boundary condition but bear on divertor asymmetry studies: \code{model600}
lacks the $\nabla B$ terms in the energy equations, and it lacks the thermoelectric heat-flux
divergence $\nabla\!\cdot\!(0.71\,T_e j_\parallel / e)$. Both appear in SOLPS and both act in the
inner/outer asymmetry channel.

\paragraph{Next steps.} With the electrical closure established, the physics questions become
accessible: extracting $\Phi(R)$ and hence $E_r$ along each target; comparing the in--out density
ratio against an otherwise identical run without the sheath condition; and the differential target
biasing scan (\code{sheath\_V\_wall\_asym}), which is now unblocked since a uniform bias is a pure
gauge and only a \emph{difference} between targets creates a potential drop across the private flux
region.

\section*{Acknowledgement of method}

The failure recorded in Section~\ref{sec:nodalfail} is the substantive lesson of this work: a
constraint can be satisfied to four decimal places in the norm one is measuring and still be wrong by
tens of percent in the quantity one cares about. The diagnostic that resolved it --- instrumenting
the weak residual as a pure observable, before changing any code --- cost far less than the
sequence of tuning experiments it made unnecessary.

\begin{thebibliography}{9}
\bibitem{stangeby} P.~C. Stangeby, \emph{The Plasma Boundary of Magnetic Fusion Devices},
  IOP Publishing, 2000. (Sheath characteristic, eq.~2.68.)
\bibitem{hoelzl2021} M.~H\"olzl \emph{et al.}, ``The JOREK non-linear extended MHD code and
  applications to large-scale instabilities and their control in magnetically confined fusion
  plasmas'', \emph{Nucl.\ Fusion} \textbf{61} 065001 (2021).
\bibitem{artola} F.~J. Artola \emph{et al.}, on sheath boundary conditions for the electrostatic
  potential in JOREK. (Saturation current, eqs.~5 and 16.)
\bibitem{chodura} R.~Chodura, ``Plasma--wall transition in an oblique magnetic field'',
  \emph{Phys.\ Fluids} \textbf{25} 1628 (1982).
\end{thebibliography}

\end{document}
